\begin{quadro}[htb]
\centering
\caption{Paradigmas de agentes segundo Müller (1999) e sua aplicação no sistema proposto}
\label{qua:muller_agentes}
\small
\begin{tabular}{|p{2.8cm}|p{2.8cm}|p{5.2cm}|p{4.2cm}|}
\hline
\textbf{Categoria} & 
\textbf{Classificação} & 
\textbf{Descrição Teórica} & 
\textbf{Aplicação no \newline Monitoramento Hídrico} \\ \hline

Arquitetura de \newline Controle & 
Reativa & 
Decisões baseadas em estímulos diretos, 
sem modelo interno do mundo. & 
Leitura em tempo real e processamento 
imediato de alertas básicos. \\ \cline{2-4}

& 
Deliberativa & 
Raciocínio lógico, planejamento e representações 
internas complexas. & 
Agente de IA interpretando as anomalias 
e formulando diagnósticos. \\ \cline{2-4}

& 
Híbrida & 
Combinação de camadas reativas (tempo real) 
e deliberativas (metas). & 
Arquitetura central proposta para o monitoramento 
de alta frequência. \\ \cline{2-4}

& 
De Interação & 
Foco na coordenação, cooperação e comunicação 
entre agentes. & 
Orquestração do Sistema Multiagente (SMA) 
para monitoramento hídrico. \\ \hline

Taxonomia: \newline Estado Material & 
Hardware & 
Agentes com presença física, interagindo 
via sensores físicos. & 
Sondas multiparamétricas instaladas 
\textit{in situ} nos rios. \\ \cline{2-4}

& 
Software & 
Agentes virtuais operando em ambientes 
computacionais ou digitais. & 
Agentes cognitivos baseados em LLMs 
e modelos de IA. \\ \hline

Taxonomia: \newline Modo de Interação & 
Autônomo & 
Foco na relação direta e independente entre o 
próprio agente e seu ambiente. & 
Sonda hídrica coletando e transmitindo dados 
de forma isolada por telemetria. \\ \cline{2-4}

& 
Multiagente & 
Presença de outros agentes no ambiente, com 
cooperação e compartilhamento de informações. & 
Nós e sondas cooperando para analisar a 
qualidade da água ao longo de toda a bacia. \\ \cline{2-4}

& 
Assistente & 
Projetados para interagir diretamente com 
humanos e agir em seu benefício. & 
Interface em linguagem natural (LLM) gerando 
relatórios claros para o gestor. \\ \hline
\end{tabular}

\vspace{0.1cm}
{\footnotesize Fonte – Adaptado de \citeonline{muller1999}}
\end{quadro}
